\documentclass[10pt,letterpaper]{article}

\usepackage[margin=0.72in]{geometry}
\usepackage{newtxtext,newtxmath}
\usepackage{booktabs}
\usepackage{amsmath}
\usepackage{graphicx}
\usepackage{array}
\usepackage{caption}
\usepackage{float}
\usepackage{microtype}
\usepackage{fancyhdr}

\captionsetup{font=small,labelfont=bf,skip=4pt}
\captionsetup[table]{position=top}
\renewcommand{\thetable}{S\arabic{table}}
\renewcommand{\thefigure}{S\arabic{figure}}
\renewcommand{\arraystretch}{1.1}
\setlength{\parskip}{0.28em}
\setlength{\parindent}{1.05em}
\setlength{\tabcolsep}{5pt}

\pagestyle{fancy}
\fancyhf{}
\renewcommand{\headrulewidth}{0.4pt}
\fancyhead[L]{\small Supplementary Material}
\fancyhead[R]{\small SPEAR-3D}
\fancyfoot[C]{\thepage}

\begin{document}
\begin{center}
{\large\bfseries Supplementary Material}\\[0.55em]
{\normalsize SPEAR-3D: From Streaming Reconstruction to Explicit Reasoning---\\
A Neuro-Symbolic Approach for 3D Visual Grounding}
\end{center}
\vspace{0.15em}
\thispagestyle{fancy}

\noindent
This file is supplementary material.
It is submitted in the zip with the demonstration video.
It is not part of the manuscript PDF.
The demonstration is the accompanying file spear3d\_demo.mp4.

\section{Additional Experiments}
The tables in this file are not repeated in the main text.
Acc@0.3m and Acc@0.5m are centroid distances in metres over queries with a finite distance, with no class match.
T-Acc requires the referred object to be the unique highest point IoU of the selected instance on the ScanNet ground-truth mesh.
No selection is incorrect.

\subsection{NS-Grounding}
\begin{table}[H]
\centering
\caption{NS-Grounding on NR3D and SR3D.
The GLM-5 row without NS on SR3D is 51.9 / 52.5.
Bold marks the highest Acc@0.5m on each dataset.}
\label{tab:module_ablation}
\setlength{\tabcolsep}{6pt}
\begin{tabular}{llccccc}
\toprule
 & & \multicolumn{5}{c}{NR3D} \\
\cmidrule(lr){3-7}
LLM & NS & MeanDist & Acc@0.3m & Acc@0.5m & T-Acc@0.25 & T-Acc@0.50 \\
\midrule
Qwen3-plus & no & 1.36 & 38.7 & 43.5 & 33.2 & 28.2 \\
Qwen3-plus & yes & 0.91 & 56.3 & \textbf{58.8} & 50.4 & 40.8 \\
deepseek-v3.2 & no & 1.33 & 42.0 & 44.6 & 30.9 & 26.0 \\
deepseek-v3.2 & yes & 1.27 & 51.1 & 54.3 & 46.0 & 39.4 \\
GLM-5 & no & 1.35 & 37.8 & 41.5 & 32.3 & 29.1 \\
GLM-5 & yes & 1.18 & 50.2 & 53.4 & 44.4 & 36.2 \\
\midrule
 & & \multicolumn{5}{c}{SR3D} \\
\cmidrule(lr){3-7}
LLM & NS & MeanDist & Acc@0.3m & Acc@0.5m & T-Acc@0.25 & T-Acc@0.50 \\
\midrule
Qwen3-plus & no & 1.32 & 53.8 & 54.8 & 37.2 & 31.1 \\
Qwen3-plus & yes & 0.96 & 64.7 & 65.9 & 61.4 & 52.6 \\
deepseek-v3.2 & no & 1.24 & 53.1 & 54.1 & 36.3 & 29.7 \\
deepseek-v3.2 & yes & 1.00 & 64.8 & 66.0 & 64.0 & 57.0 \\
GLM-5 & no & 1.28 & 51.9 & 52.5 & 35.7 & 29.5 \\
GLM-5 & yes & 0.97 & 65.6 & \textbf{67.0} & 63.2 & 53.8 \\
\bottomrule
\end{tabular}
\end{table}

\subsection{Node--Relation Weight}
\begin{table}[H]
\centering
\caption{Sensitivity of the weighted geometric mean on NR3D with Qwen3-plus.}
\label{tab:graph_balance}
\begin{tabular}{cccccc}
\toprule
\(\lambda_u\) & MeanDist & Acc@0.3m & Acc@0.5m & T-Acc@0.25 & T-Acc@0.50 \\
\midrule
0.25 & 1.2550 & 48.1 & 51.6 & -- & -- \\
0.50 & 0.9140 & 56.3 & 58.8 & 50.4 & 40.8 \\
0.75 & 1.0840 & 53.4 & 57.8 & -- & -- \\
\bottomrule
\end{tabular}
\end{table}

\subsection{End-to-End Frame Rate}
\begin{table}[H]
\centering
\caption{End-to-end FPS on an RTX 3090 with keyframe interval 10.
A is SAM3 with label binding.
C is CropFormer with per-crop CLIP.}
\label{tab:fps}
\begin{tabular}{lcc}
\toprule
Scene & A & C \\
\midrule
scene0432\_00 & 11.42 & 5.50 \\
scene0527\_00 & 10.71 & 4.82 \\
scene0559\_00 & 10.62 & 4.41 \\
scene0494\_00 & 10.82 & 6.60 \\
scene0689\_00 & 11.04 & 5.25 \\
Mean & 10.92 & 5.32 \\
\bottomrule
\end{tabular}
\end{table}

\subsection{Per-Stage Time}
\begin{table}[H]
\centering
\caption{Per-stage wall time.
Rates are reciprocals of those times.}
\label{tab:stages}
\begin{tabular}{lcc}
\toprule
Stage & A & C \\
\midrule
Segmentation & 241 ms (4.16/s) & 1099 ms (0.91/s) \\
Fusion & 330 ms (3.03/s) & 453 ms (2.21/s) \\
Merge & 487 ms (2.05/s) & 1307 ms (0.77/s) \\
\bottomrule
\end{tabular}
\end{table}

\subsection{Prompt Vocabulary}
\begin{table}[H]
\centering
\caption{Segmentation latency as the prompt vocabulary grows.}
\label{tab:vocab}
\begin{tabular}{lcccc}
\toprule
 & \multicolumn{4}{c}{Vocabulary size} \\
\cmidrule(lr){2-5}
 & 10 & 30 & 50 & 100 \\
\midrule
Latency (ms) & 311 & 509 & 762 & 1627 \\
Seg.\ rate & 3.22 & 1.97 & 1.31 & 0.61 \\
Peak MB & 16105 & 16098 & 16095 & 16084 \\
Fragments & 8.3 & 11.3 & 12.1 & 16.8 \\
\bottomrule
\end{tabular}
\end{table}

\noindent
Segmentation takes 241~ms (4.16/s) against 1099~ms (0.91/s), fusion 330~ms (3.03/s) against 453~ms (2.21/s), and merge 487~ms (2.05/s) against 1307~ms (0.77/s).
For vocabulary sizes 10, 30, 50, and 100, latency is 311, 509, 762, and 1627~ms, segmentation rate is 3.22, 1.97, 1.31, and 0.61, peak memory is 16105, 16098, 16095, and 16084~MB, and the mean fragment count is 8.3, 11.3, 12.1, and 16.8.
About 10~FPS is the end-to-end mean 10.92 of arm A.
The published OnlineAnySeg figure of about 15~FPS is on an RTX 4090.

\subsection{SceneNN Coverage}
\begin{table}[H]
\centering
\caption{SceneNN prompt-vocabulary coverage.
Out-of-vocabulary class-matched coverage is zero because that class is not a prompt.}
\label{tab:scenenn}
\begin{tabular}{lc}
\toprule
Quantity & Value \\
\midrule
In-vocabulary instances & 61.3\% \\
Out-of-vocabulary instances & 38.7\% \\
\texttt{otherprops} share & 33.8\% \\
In-vocab. recall @0.50 / @0.25 & 56.7\% / 85.1\% \\
Out-of-vocab. recall @0.50 / @0.25 & 37.8\% / 61.3\% \\
Class-matched coverage, in / out & 61.1\% / 0.0\% \\
Unlabeled vertices & 28.2\% \\
\bottomrule
\end{tabular}
\end{table}

\subsection{Label Binding}
\begin{table}[H]
\centering
\caption{Label binding and post-hoc CLIP on the same SAM3 masks.
The manuscript cites label-binding mIoU 30.6 against 15.3 for post-hoc CLIP.}
\label{tab:bind}
\begin{tabular}{lccc}
\toprule
Method & mIoU & mAcc & Queries \\
\midrule
Label binding & 30.6 & 39.8 & 0 \\
Post-hoc CLIP argmax & 15.3 & 26.0 & 1 \\
\bottomrule
\end{tabular}
\end{table}

\subsection{Cited OVI-MAP Scores}
\begin{table}[H]
\centering
\caption{Percent scores beside OVI-MAP Table~3.
The OVI-MAP row is cited and was not re-run.
The rows are not a paired comparison.}
\label{tab:ovimap}
\begin{tabular}{lcc}
\toprule
Method & mIoU & mAcc \\
\midrule
OVI-MAP & 17.5 & 27.6 \\
SPEAR-3D & 30.6 & 39.8 \\
\bottomrule
\end{tabular}
\end{table}

\subsection{Synthetic Degradation}
\begin{table}[H]
\centering
\caption{Class-agnostic AP under blur, depth noise, and frame drop.
Blur and noise are not monotone.
Frame drop is monotone.}
\label{tab:robust}
\begin{tabular}{lccc}
\toprule
Condition & AP & AP50 & AP25 \\
\midrule
Clean & 22.6 & 40.4 & 60.4 \\
Blur 7 px & 21.2 & 38.5 & 60.9 \\
Blur 21 px & 22.6 & 42.0 & 63.2 \\
Blur 41 px & 21.7 & 40.7 & 59.7 \\
Depth \(b{=}0.005\), holes 3\% & 21.9 & 40.3 & 63.2 \\
Depth \(b{=}0.02\), holes 10\% & 22.0 & 39.9 & 63.4 \\
Depth \(b{=}0.04\), holes 20\% & 21.5 & 38.6 & 61.7 \\
Keep 133 frames & 19.3 & 34.8 & 54.4 \\
Keep 100 frames & 17.0 & 32.6 & 50.0 \\
Keep 67 frames & 13.2 & 27.6 & 42.9 \\
\bottomrule
\end{tabular}
\end{table}

\noindent
At 4~m, \(\sigma\) is about 0.54, 0.66, and 0.82~mm.
The Gaussian term is below one millimetre.
The main change along the depth rows is the hole rate.

\subsection{Segmenter and Binding}
\begin{table}[H]
\centering
\caption{Differences of the rounded \(2\times 2\) cells, in percentage points.}
\label{tab:recondiff}
\begin{tabular}{llccc}
\toprule
Contrast & Meaning & \(\Delta\)AP & \(\Delta\)AP50 & \(\Delta\)AP25 \\
\midrule
A\(-\)B & Binding, SAM3 fixed & +2.4 & +3.0 & +1.8 \\
D\(-\)C & Binding, CropFormer fixed & +3.2 & +3.7 & +4.4 \\
A\(-\)D & Backbone, binding fixed & +0.7 & +0.6 & +2.6 \\
B\(-\)C & Backbone, per-crop fixed & +1.5 & +1.3 & +5.2 \\
A\(-\)C & Full method vs.\ OnlineAnySeg & +3.9 & +4.3 & +7.0 \\
\bottomrule
\end{tabular}
\end{table}

\noindent
On AP and AP50, label binding is the larger effect when the segmenter is fixed.
On AP25, the backbone gap under label binding (\(+2.6\)) exceeds the binding gap on SAM3 (\(+1.8\)).

\section{Additional Figures}
These two figures use the same chair example as Figure 4 of the manuscript.

\begin{figure}[H]
\centering
\includegraphics[width=0.86\linewidth]{scene_graph.pdf}
\caption{3D scene graph topology.
The graph encodes explicit spatial relations as the substrate for structural graph matching.}
\label{fig:s_scene}
\end{figure}

\begin{figure}[H]
\centering
\includegraphics[width=0.86\linewidth]{pred_3d.pdf}
\caption{Final 3D spatial localization.
The red bounding box indicates the predicted target extent in the reconstructed scene.}
\label{fig:s_pred}
\end{figure}

\section{Robot Demonstration Video}
\label{sec:robotvideo}
This revision does not add a quantitative physical-robot experiment.
The demonstration is the accompanying file spear3d\_demo.mp4.

\end{document}
